%%%%%%%%%%%%%%%%%%%%%%%%%%%%%%%%%%%%%%%%%
% The Legrand Orange Book
% LaTeX Template
% Version 2.4 (26/09/2018)
%
% This template was downloaded from:
% http://www.LaTeXTemplates.com
%
% Original author:
% Mathias Legrand (legrand.mathias@gmail.com) with modifications by:
% Vel (vel@latextemplates.com)
%
% License:
% CC BY-NC-SA 3.0 (http://creativecommons.org/licenses/by-nc-sa/3.0/)
%
% Compiling this template:
% This template uses biber for its bibliography and makeindex for its index.
% When you first open the template, compile it from the command line with the 
% commands below to make sure your LaTeX distribution is configured correctly:
%
% 1) pdflatex main
% 2) makeindex main.idx -s StyleInd.ist
% 3) biber main
% 4) pdflatex main x 2
%
% After this, when you wish to update the bibliography/index use the appropriate
% command above and make sure to compile with pdflatex several times 
% afterwards to propagate your changes to the document.
%
% This template also uses a number of packages which may need to be
% updated to the newest versions for the template to compile. It is strongly
% recommended you update your LaTeX distribution if you have any
% compilation errors.
%
% Important note:
% Chapter heading images should have a 2:1 width:height ratio,
% e.g. 920px width and 460px height.
%
%%%%%%%%%%%%%%%%%%%%%%%%%%%%%%%%%%%%%%%%%

%----------------------------------------------------------------------------------------
%	PACKAGES AND OTHER DOCUMENT CONFIGURATIONS
%----------------------------------------------------------------------------------------

\documentclass[11pt,fleqn]{book} % Default font size and left-justified equations

%%%%%%%%%%%%%%%%%%%%%%%%%%%%%%%%%%%%%%%%%
% The Legrand Orange Book
% Structural Definitions File
% Version 2.1 (26/09/2018)
%
% Original author:
% Mathias Legrand (legrand.mathias@gmail.com) with modifications by:
% Vel (vel@latextemplates.com)
% 
% This file was downloaded from:
% http://www.LaTeXTemplates.com
%
% License:
% CC BY-NC-SA 3.0 (http://creativecommons.org/licenses/by-nc-sa/3.0/)
%
%%%%%%%%%%%%%%%%%%%%%%%%%%%%%%%%%%%%%%%%%

%----------------------------------------------------------------------------------------
%	VARIOUS REQUIRED PACKAGES AND CONFIGURATIONS
%----------------------------------------------------------------------------------------

\usepackage{graphicx} % Required for including pictures
\graphicspath{{Pictures/}} % Specifies the directory where pictures are stored

\usepackage{lipsum} % Inserts dummy text

\usepackage{tikz} % Required for drawing custom shapes

%\usepackage[english]{babel} % English language/hyphenation
%\usepackage[russian]{babel}

\usepackage[T2A]{fontenc}
\usepackage[utf8]{inputenc}
\usepackage[russian]{babel}


\usepackage{wrapfig}
\usepackage{mathtools}
\usepackage{ stmaryrd }

\usepackage{enumitem} % Customize lists
\setlist{nolistsep} % Reduce spacing between bullet points and numbered lists

\usepackage{booktabs} % Required for nicer horizontal rules in tables

\usepackage{xcolor} % Required for specifying colors by name
\definecolor{ocre}{RGB}{243,102,25} % Define the orange color used for highlighting throughout the book

%----------------------------------------------------------------------------------------
%	MARGINS
%----------------------------------------------------------------------------------------

\usepackage{geometry} % Required for adjusting page dimensions and margins

\geometry{
	paper=a4paper, % Paper size, change to letterpaper for US letter size
	top=3cm, % Top margin
	bottom=3cm, % Bottom margin
	left=3cm, % Left margin
	right=3cm, % Right margin
	headheight=14pt, % Header height
	footskip=1.4cm, % Space from the bottom margin to the baseline of the footer
	headsep=10pt, % Space from the top margin to the baseline of the header
	%showframe, % Uncomment to show how the type block is set on the page
}

%----------------------------------------------------------------------------------------
%	FONTS
%----------------------------------------------------------------------------------------
%\usepackage{XCharter}
%\usepackage[noae]{Sweave}
%\usepackage{gentium}
%\usepackage{avant} % Use the Avantgarde font for headings
%\usepackage{times} % Use the Times font for headings
%\usepackage{mathptmx} % Use the Adobe Times Roman as the default text font together with math symbols from the Sym­bol, Chancery and Com­puter Modern fonts


\usepackage{microtype} % Slightly tweak font spacing for aesthetics
\usepackage[utf8]{inputenc} % Required for including letters with accents
\usepackage[T1]{fontenc} % Use 8-bit encoding that has 256 glyphs


%----------------------------------------------------------------------------------------
%	BIBLIOGRAPHY AND INDEX
%----------------------------------------------------------------------------------------

\usepackage[style=numeric,citestyle=numeric,sorting=nyt,sortcites=true,autopunct=true,babel=hyphen,hyperref=true,abbreviate=false,backref=true,backend=biber]{biblatex}
\addbibresource{bibliography.bib} % BibTeX bibliography file
\defbibheading{bibempty}{}

\usepackage{calc} % For simpler calculation - used for spacing the index letter headings correctly
\usepackage{makeidx} % Required to make an index
\makeindex % Tells LaTeX to create the files required for indexing

%----------------------------------------------------------------------------------------
%	MAIN TABLE OF CONTENTS
%----------------------------------------------------------------------------------------

\usepackage{titletoc} % Required for manipulating the table of contents

\contentsmargin{0cm} % Removes the default margin

% Part text styling (this is mostly taken care of in the PART HEADINGS section of this file)
\titlecontents{part}
	[0cm] % Left indentation
	{\addvspace{20pt}\bfseries} % Spacing and font options for parts
	{}
	{}
	{}

% Chapter text styling
\titlecontents{chapter}
	[1.25cm] % Left indentation
	{\addvspace{12pt}\large\sffamily\bfseries} % Spacing and font options for chapters
	{\color{ocre!60}\contentslabel[\Large\thecontentslabel]{1.25cm}\color{ocre}} % Formatting of numbered sections of this type
	{\color{ocre}} % Formatting of numberless sections of this type
	{\color{ocre!60}\normalsize\;\titlerule*[.5pc]{.}\;\thecontentspage} % Formatting of the filler to the right of the heading and the page number

% Section text styling
\titlecontents{section}
	[1.25cm] % Left indentation
	{\addvspace{3pt}\sffamily\bfseries} % Spacing and font options for sections
	{\contentslabel[\thecontentslabel]{1.25cm}} % Formatting of numbered sections of this type
	{} % Formatting of numberless sections of this type
	{\hfill\color{black}\thecontentspage} % Formatting of the filler to the right of the heading and the page number

% Subsection text styling
\titlecontents{subsection}
	[1.25cm] % Left indentation
	{\addvspace{1pt}\sffamily\small} % Spacing and font options for subsections
	{\contentslabel[\thecontentslabel]{1.25cm}} % Formatting of numbered sections of this type
	{} % Formatting of numberless sections of this type
	{\ \titlerule*[.5pc]{.}\;\thecontentspage} % Formatting of the filler to the right of the heading and the page number

% Figure text styling
\titlecontents{figure}
	[1.25cm] % Left indentation
	{\addvspace{1pt}\sffamily\small} % Spacing and font options for figures
	{\thecontentslabel\hspace*{1em}} % Formatting of numbered sections of this type
	{} % Formatting of numberless sections of this type
	{\ \titlerule*[.5pc]{.}\;\thecontentspage} % Formatting of the filler to the right of the heading and the page number

% Table text styling
\titlecontents{table}
	[1.25cm] % Left indentation
	{\addvspace{1pt}\sffamily\small} % Spacing and font options for tables
	{\thecontentslabel\hspace*{1em}} % Formatting of numbered sections of this type
	{} % Formatting of numberless sections of this type
	{\ \titlerule*[.5pc]{.}\;\thecontentspage} % Formatting of the filler to the right of the heading and the page number

%----------------------------------------------------------------------------------------
%	MINI TABLE OF CONTENTS IN PART HEADS
%----------------------------------------------------------------------------------------

% Chapter text styling
\titlecontents{lchapter}
	[0em] % Left indentation
	{\addvspace{15pt}\large\sffamily\bfseries} % Spacing and font options for chapters
	{\color{ocre}\contentslabel[\Large\thecontentslabel]{1.25cm}\color{ocre}} % Chapter number
	{}  
	{\color{ocre}\normalsize\sffamily\bfseries\;\titlerule*[.5pc]{.}\;\thecontentspage} % Page number

% Section text styling
\titlecontents{lsection}
	[0em] % Left indentation
	{\sffamily\small} % Spacing and font options for sections
	{\contentslabel[\thecontentslabel]{1.25cm}} % Section number
	{}
	{}

% Subsection text styling (note these aren't shown by default, display them by searchings this file for tocdepth and reading the commented text)
\titlecontents{lsubsection}
	[.5em] % Left indentation
	{\sffamily\footnotesize} % Spacing and font options for subsections
	{\contentslabel[\thecontentslabel]{1.25cm}}
	{}
	{}

%----------------------------------------------------------------------------------------
%	HEADERS AND FOOTERS
%----------------------------------------------------------------------------------------

\usepackage{fancyhdr} % Required for header and footer configuration

\pagestyle{fancy} % Enable the custom headers and footers

\renewcommand{\chaptermark}[1]{\markboth{\sffamily\normalsize\bfseries\chaptername\ \thechapter.\ #1}{}} % Styling for the current chapter in the header
\renewcommand{\sectionmark}[1]{\markright{\sffamily\normalsize\thesection\hspace{5pt}#1}{}} % Styling for the current section in the header

\fancyhf{} % Clear default headers and footers
\fancyhead[LE,RO]{\sffamily\normalsize\thepage} % Styling for the page number in the header
\fancyhead[LO]{\rightmark} % Print the nearest section name on the left side of odd pages
\fancyhead[RE]{\leftmark} % Print the current chapter name on the right side of even pages
%\fancyfoot[C]{\thepage} % Uncomment to include a footer

\renewcommand{\headrulewidth}{0.5pt} % Thickness of the rule under the header

\fancypagestyle{plain}{% Style for when a plain pagestyle is specified
	\fancyhead{}\renewcommand{\headrulewidth}{0pt}%
}

% Removes the header from odd empty pages at the end of chapters
\makeatletter
\renewcommand{\cleardoublepage}{
\clearpage\ifodd\c@page\else
\hbox{}
\vspace*{\fill}
\thispagestyle{empty}
\newpage
\fi}

%----------------------------------------------------------------------------------------
%	THEOREM STYLES
%----------------------------------------------------------------------------------------

\usepackage{amsmath,amsfonts,amssymb,amsthm} % For math equations, theorems, symbols, etc

\newcommand{\intoo}[2]{\mathopen{]}#1\,;#2\mathclose{[}}
\newcommand{\ud}{\mathop{\mathrm{{}d}}\mathopen{}}
\newcommand{\intff}[2]{\mathopen{[}#1\,;#2\mathclose{]}}
\renewcommand{\qedsymbol}{$\blacksquare$}
\newtheorem{notation}{Notation}[chapter]

% Boxed/framed environments
\newtheoremstyle{ocrenumbox}% Theorem style name
{0pt}% Space above
{0pt}% Space below
{\normalfont}% Body font
{}% Indent amount
{\small\bf\sffamily\color{ocre}}% Theorem head font
{\;}% Punctuation after theorem head
{0.25em}% Space after theorem head
{\small\sffamily\color{ocre}\thmname{#1}\nobreakspace\thmnumber{\@ifnotempty{#1}{}\@upn{#2}}% Theorem text (e.g. Theorem 2.1)
\thmnote{\nobreakspace\the\thm@notefont\sffamily\bfseries\color{black}---\nobreakspace#3.}} % Optional theorem note

\newtheoremstyle{blacknumex}% Theorem style name
{5pt}% Space above
{5pt}% Space below
{\normalfont}% Body font
{} % Indent amount
{\small\bf\sffamily}% Theorem head font
{\;}% Punctuation after theorem head
{0.25em}% Space after theorem head
{\small\sffamily{\tiny\ensuremath{\blacksquare}}\nobreakspace\thmname{#1}\nobreakspace\thmnumber{\@ifnotempty{#1}{}\@upn{#2}}% Theorem text (e.g. Theorem 2.1)
\thmnote{\nobreakspace\the\thm@notefont\sffamily\bfseries---\nobreakspace#3.}}% Optional theorem note

\newtheoremstyle{blacknumbox} % Theorem style name
{0pt}% Space above
{0pt}% Space below
{\normalfont}% Body font
{}% Indent amount
{\small\bf\sffamily}% Theorem head font
{\;}% Punctuation after theorem head
{0.25em}% Space after theorem head
{\small\sffamily\thmname{#1}\nobreakspace\thmnumber{\@ifnotempty{#1}{}\@upn{#2}}% Theorem text (e.g. Theorem 2.1)
\thmnote{\nobreakspace\the\thm@notefont\sffamily\bfseries---\nobreakspace#3.}}% Optional theorem note

% Non-boxed/non-framed environments
\newtheoremstyle{ocrenum}% Theorem style name
{5pt}% Space above
{5pt}% Space below
{\normalfont}% Body font
{}% Indent amount
{\small\bf\sffamily\color{ocre}}% Theorem head font
{\;}% Punctuation after theorem head
{0.25em}% Space after theorem head
{\small\sffamily\color{ocre}\thmname{#1}\nobreakspace\thmnumber{\@ifnotempty{#1}{}\@upn{#2}}% Theorem text (e.g. Theorem 2.1)
\thmnote{\nobreakspace\the\thm@notefont\sffamily\bfseries\color{black}---\nobreakspace#3.}} % Optional theorem note
\makeatother

% Defines the theorem text style for each type of theorem to one of the three styles above
\newcounter{dummy} 
\numberwithin{dummy}{section}
\theoremstyle{ocrenumbox}
\newtheorem{theoremeT}[dummy]{Теорема}
\newtheorem{problem}{Задача}[chapter]
\newtheorem{exerciseT}{Упражнение}[chapter]
\theoremstyle{blacknumex}
\newtheorem{exampleT}{Пример}[chapter]
\theoremstyle{blacknumbox}
\newtheorem{vocabulary}{Словарь}[chapter]
\newtheorem{definitionT}{Определение}[section]
\newtheorem{corollaryT}[dummy]{Следствие}
\theoremstyle{ocrenum}
\newtheorem{proposition}[dummy]{Предложение}

%----------------------------------------------------------------------------------------
%	DEFINITION OF COLORED BOXES
%----------------------------------------------------------------------------------------

\RequirePackage[framemethod=default]{mdframed} % Required for creating the theorem, definition, exercise and corollary boxes

% Theorem box
\newmdenv[skipabove=7pt,
skipbelow=7pt,
backgroundcolor=black!5,
linecolor=ocre,
innerleftmargin=5pt,
innerrightmargin=5pt,
innertopmargin=5pt,
leftmargin=0cm,
rightmargin=0cm,
innerbottommargin=5pt]{tBox}

% Exercise box	  
\newmdenv[skipabove=7pt,
skipbelow=7pt,
rightline=false,
leftline=true,
topline=false,
bottomline=false,
backgroundcolor=ocre!10,
linecolor=ocre,
innerleftmargin=5pt,
innerrightmargin=5pt,
innertopmargin=5pt,
innerbottommargin=5pt,
leftmargin=0cm,
rightmargin=0cm,
linewidth=4pt]{eBox}	

% Definition box
\newmdenv[skipabove=7pt,
skipbelow=7pt,
rightline=false,
leftline=true,
topline=false,
bottomline=false,
linecolor=ocre,
innerleftmargin=5pt,
innerrightmargin=5pt,
innertopmargin=0pt,
leftmargin=0cm,
rightmargin=0cm,
linewidth=4pt,
innerbottommargin=0pt]{dBox}	

% Corollary box
\newmdenv[skipabove=7pt,
skipbelow=7pt,
rightline=false,
leftline=true,
topline=false,
bottomline=false,
linecolor=gray,
backgroundcolor=black!5,
innerleftmargin=5pt,
innerrightmargin=5pt,
innertopmargin=5pt,
leftmargin=0cm,
rightmargin=0cm,
linewidth=4pt,
innerbottommargin=5pt]{cBox}

% Creates an environment for each type of theorem and assigns it a theorem text style from the "Theorem Styles" section above and a colored box from above
\newenvironment{theorem}{\begin{tBox}\begin{theoremeT}}{\end{theoremeT}\end{tBox}}
\newenvironment{exercise}{\begin{eBox}\begin{exerciseT}}{\hfill{\color{ocre}\tiny\ensuremath{\blacksquare}}\end{exerciseT}\end{eBox}}				  
\newenvironment{definition}{\begin{dBox}\begin{definitionT}}{\end{definitionT}\end{dBox}}	
\newenvironment{example}{\begin{exampleT}}{\hfill{\tiny\ensuremath{\blacksquare}}\end{exampleT}}		
\newenvironment{corollary}{\begin{cBox}\begin{corollaryT}}{\end{corollaryT}\end{cBox}}	

%----------------------------------------------------------------------------------------
%	REMARK ENVIRONMENT
%----------------------------------------------------------------------------------------

\newenvironment{remark}{\par\vspace{10pt}\small % Vertical white space above the remark and smaller font size
\begin{list}{}{
\leftmargin=35pt % Indentation on the left
\rightmargin=25pt}\item\ignorespaces % Indentation on the right
\makebox[-2.5pt]{\begin{tikzpicture}[overlay]
\node[draw=ocre!60,line width=1pt,circle,fill=ocre!25,font=\sffamily\bfseries,inner sep=2pt,outer sep=0pt] at (-15pt,0pt){\textcolor{ocre}{!}};\end{tikzpicture}} % Orange R in a circle
\advance\baselineskip -1pt}{\end{list}\vskip5pt} % Tighter line spacing and white space after remark

%----------------------------------------------------------------------------------------
%	SECTION NUMBERING IN THE MARGIN
%----------------------------------------------------------------------------------------

\makeatletter
\renewcommand{\@seccntformat}[1]{\llap{\textcolor{ocre}{\csname the#1\endcsname}\hspace{1em}}}                    
\renewcommand{\section}{\@startsection{section}{1}{\z@}
{-4ex \@plus -1ex \@minus -.4ex}
{1ex \@plus.2ex }
{\normalfont\large\sffamily\bfseries}}
\renewcommand{\subsection}{\@startsection {subsection}{2}{\z@}
{-3ex \@plus -0.1ex \@minus -.4ex}
{0.5ex \@plus.2ex }
{\normalfont\sffamily\bfseries}}
\renewcommand{\subsubsection}{\@startsection {subsubsection}{3}{\z@}
{-2ex \@plus -0.1ex \@minus -.2ex}
{.2ex \@plus.2ex }
{\normalfont\small\sffamily\bfseries}}                        
\renewcommand\paragraph{\@startsection{paragraph}{4}{\z@}
{-2ex \@plus-.2ex \@minus .2ex}
{.1ex}
{\normalfont\small\sffamily\bfseries}}

%----------------------------------------------------------------------------------------
%	PART HEADINGS
%----------------------------------------------------------------------------------------

% Numbered part in the table of contents
\newcommand{\@mypartnumtocformat}[2]{%
	\setlength\fboxsep{0pt}%
	\noindent\colorbox{ocre!20}{\strut\parbox[c][.7cm]{\ecart}{\color{ocre!70}\Large\sffamily\bfseries\centering#1}}\hskip\esp\colorbox{ocre!40}{\strut\parbox[c][.7cm]{\linewidth-\ecart-\esp}{\Large\sffamily\centering#2}}%
}

% Unnumbered part in the table of contents
\newcommand{\@myparttocformat}[1]{%
	\setlength\fboxsep{0pt}%
	\noindent\colorbox{ocre!40}{\strut\parbox[c][.7cm]{\linewidth}{\Large\sffamily\centering#1}}%
}

\newlength\esp
\setlength\esp{4pt}
\newlength\ecart
\setlength\ecart{1.2cm-\esp}
\newcommand{\thepartimage}{}%
\newcommand{\partimage}[1]{\renewcommand{\thepartimage}{#1}}%
\def\@part[#1]#2{%
\ifnum \c@secnumdepth >-2\relax%
\refstepcounter{part}%
\addcontentsline{toc}{part}{\texorpdfstring{\protect\@mypartnumtocformat{\thepart}{#1}}{\partname~\thepart\ ---\ #1}}
\else%
\addcontentsline{toc}{part}{\texorpdfstring{\protect\@myparttocformat{#1}}{#1}}%
\fi%
\startcontents%
\markboth{}{}%
{\thispagestyle{empty}%
\begin{tikzpicture}[remember picture,overlay]%
\node at (current page.north west){\begin{tikzpicture}[remember picture,overlay]%	
\fill[ocre!20](0cm,0cm) rectangle (\paperwidth,-\paperheight);
\node[anchor=north] at (4cm,-3.25cm){\color{ocre!40}\fontsize{220}{100}\sffamily\bfseries\thepart}; 
\node[anchor=south east] at (\paperwidth-1cm,-\paperheight+1cm){\parbox[t][][t]{8.5cm}{
\printcontents{l}{0}{\setcounter{tocdepth}{1}}% The depth to which the Part mini table of contents displays headings; 0 for chapters only, 1 for chapters and sections and 2 for chapters, sections and subsections
}};
\node[anchor=north east] at (\paperwidth-1.5cm,-3.25cm){\parbox[t][][t]{15cm}{\strut\raggedleft\color{white}\fontsize{30}{30}\sffamily\bfseries#2}};
\end{tikzpicture}};
\end{tikzpicture}}%
\@endpart}
\def\@spart#1{%
\startcontents%
\phantomsection
{\thispagestyle{empty}%
\begin{tikzpicture}[remember picture,overlay]%
\node at (current page.north west){\begin{tikzpicture}[remember picture,overlay]%	
\fill[ocre!20](0cm,0cm) rectangle (\paperwidth,-\paperheight);
\node[anchor=north east] at (\paperwidth-1.5cm,-3.25cm){\parbox[t][][t]{15cm}{\strut\raggedleft\color{white}\fontsize{30}{30}\sffamily\bfseries#1}};
\end{tikzpicture}};
\end{tikzpicture}}
\addcontentsline{toc}{part}{\texorpdfstring{%
\setlength\fboxsep{0pt}%
\noindent\protect\colorbox{ocre!40}{\strut\protect\parbox[c][.7cm]{\linewidth}{\Large\sffamily\protect\centering #1\quad\mbox{}}}}{#1}}%
\@endpart}
\def\@endpart{\vfil\newpage
\if@twoside
\if@openright
\null
\thispagestyle{empty}%
\newpage
\fi
\fi
\if@tempswa
\twocolumn
\fi}

%----------------------------------------------------------------------------------------
%	CHAPTER HEADINGS
%----------------------------------------------------------------------------------------

% A switch to conditionally include a picture, implemented by Christian Hupfer
\newif\ifusechapterimage
\usechapterimagetrue
\newcommand{\thechapterimage}{}%
\newcommand{\chapterimage}[1]{\ifusechapterimage\renewcommand{\thechapterimage}{#1}\fi}%
\newcommand{\autodot}{.}
\def\@makechapterhead#1{%
{\parindent \z@ \raggedright \normalfont
\ifnum \c@secnumdepth >\m@ne
\if@mainmatter
\begin{tikzpicture}[remember picture,overlay]
\node at (current page.north west)
{\begin{tikzpicture}[remember picture,overlay]
\node[anchor=north west,inner sep=0pt] at (0,0) {\ifusechapterimage\includegraphics[width=\paperwidth]{\thechapterimage}\fi};
\draw[anchor=west] (\Gm@lmargin,-9cm) node [line width=2pt,rounded corners=15pt,draw=ocre,fill=white,fill opacity=0.5,inner sep=15pt]{\strut\makebox[22cm]{}};
\draw[anchor=west] (\Gm@lmargin+.3cm,-9cm) node {\huge\sffamily\bfseries\color{black}\thechapter\autodot~#1\strut};
\end{tikzpicture}};
\end{tikzpicture}
\else
\begin{tikzpicture}[remember picture,overlay]
\node at (current page.north west)
{\begin{tikzpicture}[remember picture,overlay]
\node[anchor=north west,inner sep=0pt] at (0,0) {\ifusechapterimage\includegraphics[width=\paperwidth]{\thechapterimage}\fi};
\draw[anchor=west] (\Gm@lmargin,-9cm) node [line width=2pt,rounded corners=15pt,draw=ocre,fill=white,fill opacity=0.5,inner sep=15pt]{\strut\makebox[22cm]{}};
\draw[anchor=west] (\Gm@lmargin+.3cm,-9cm) node {\huge\sffamily\bfseries\color{black}#1\strut};
\end{tikzpicture}};
\end{tikzpicture}
\fi\fi\par\vspace*{270\p@}}}

%-------------------------------------------

\def\@makeschapterhead#1{%
\begin{tikzpicture}[remember picture,overlay]
\node at (current page.north west)
{\begin{tikzpicture}[remember picture,overlay]
\node[anchor=north west,inner sep=0pt] at (0,0) {\ifusechapterimage\includegraphics[width=\paperwidth]{\thechapterimage}\fi};
\draw[anchor=west] (\Gm@lmargin,-9cm) node [line width=2pt,rounded corners=15pt,draw=ocre,fill=white,fill opacity=0.5,inner sep=15pt]{\strut\makebox[22cm]{}};
\draw[anchor=west] (\Gm@lmargin+.3cm,-9cm) node {\huge\sffamily\bfseries\color{black}#1\strut};
\end{tikzpicture}};
\end{tikzpicture}
\par\vspace*{270\p@}}
\makeatother

%----------------------------------------------------------------------------------------
%	LINKS
%----------------------------------------------------------------------------------------

\usepackage{hyperref}
\hypersetup{hidelinks,backref=true,pagebackref=true,hyperindex=true,colorlinks=false,breaklinks=true,urlcolor=ocre,bookmarks=true,bookmarksopen=false}

\usepackage{bookmark}
\bookmarksetup{
open,
numbered,
addtohook={%
\ifnum\bookmarkget{level}=0 % chapter
\bookmarksetup{bold}%
\fi
\ifnum\bookmarkget{level}=-1 % part
\bookmarksetup{color=ocre,bold}%
\fi
}
}
 % Insert the commands.tex file which contains the majority of the structure behind the template

%\hypersetup{pdftitle={Title},pdfauthor={Author}} % Uncomment and fill out to include PDF metadata for the author and title of the book

%----------------------------------------------------------------------------------------

\begin{document}

%----------------------------------------------------------------------------------------
%	TITLE PAGE
%----------------------------------------------------------------------------------------

\begingroup
\thispagestyle{empty} % Suppress headers and footers on the title page
\begin{tikzpicture}[remember picture,overlay]
\node[inner sep=0pt] (background) at (current page.center) {\includegraphics[width=\paperwidth]{background.pdf}};
\draw (current page.center) node [fill=ocre!30!white,fill opacity=0.6,text opacity=1,inner sep=1cm]{\Huge\centering\parbox[c][][t]{\paperwidth}{\centering \textbf{Математика: общий курс}\\[10pt] % Book title
\textit{Дейвид Капаца}}}; % Author name
\end{tikzpicture}
\vfill
\endgroup

%----------------------------------------------------------------------------------------
%	COPYRIGHT PAGE
%----------------------------------------------------------------------------------------

\newpage
~\vfill
\thispagestyle{empty}

\noindent Copyright 2020 David Kapatsa\\ % Copyright notice

%\noindent \textsc{Published by Publisher}\\ % Publisher

%\noindent \textsc{book-website.com}\\ % URL

\noindent Licensed under the Creative Commons Attribution-NonCommercial 3.0 Unported License (the ``License''). You may not use this file except in compliance with the License. You may obtain a copy of the License at \url{http://creativecommons.org/licenses/by-nc/3.0}. Unless required by applicable law or agreed to in writing, software distributed under the License is distributed on an \textsc{``as is'' basis, without warranties or conditions of any kind}, either express or implied. See the License for the specific language governing permissions and limitations under the License.\\ % License information, replace this with your own license (if any)

\noindent \textit{First printing, July 2020} % Printing/edition date

%----------------------------------------------------------------------------------------
%	TABLE OF CONTENTS
%----------------------------------------------------------------------------------------

%\usechapterimagefalse % If you don't want to include a chapter image, use this to toggle images off - it can be enabled later with \usechapterimagetrue

\chapterimage{chapter_head_1.pdf} % Table of contents heading image


\pagestyle{empty} % Disable headers and footers for the following pages

\tableofcontents % Print the table of contents itself

%\cleardoublepage % Forces the first chapter to start on an odd page so it's on the right side of the book

\pagestyle{fancy} % Enable headers and footers again

\part{Введение}

\begin{quote}
Один из учеников Евклида спросил у него: 
\end{quote}

\chapter{(Очень) Краткая история математики}



%\part{Линейные уравнения}

\part{Системы уравнений}

\chapter{Системы уравнений. Начало.}

\section{Приводящие к системам уравнений задачи}

Рассмотрим две задачи:

\begin{problem}
25 августа 2020 года на Московской фондовой бирже ваш коллега-трейдер купил $500$ акций Русала и $100$ акций Алроса за $22700$ рублей. Зная, что одна акция Алроса стоит $67$ рублей, скажите, сколько стоит одна акция Русала?
\end{problem}

\begin{problem}
25 августа 2020 года на Московской фондовой бирже ваш коллега-трейдер купил $500$ акций Русала и $100$ акций Алроса за $22\,700$ рублей. В то же время другой ваш коллега купил $300$ акций Русала и $200$ акций Алроса за $13\,000$ рублей. Назовите цены акций Русала и Алроса.
\end{problem}
~\\~
Первая задача не требует долгих размышлений и решается без помощи систем уравнений --- решите её. Очевидно, что со второй задачей уже не всё так просто.

Покажем, как будет выглядеть решение первой задачи, если решать её новым способом.
~\\ \textit{Решение:} ~\\
Пусть $x$ рублей --- цена акции Русала, а $y$ рублей --- цена акции Алроса. Давайте запишем условие таким образом:
$$
\begin{cases}
500x + 100y = 22700\\
x = 67
\end{cases}
$$
Видим два уравнения, объединённых фигурной скобкой. Такая запись называется \textbf{системой уравнений}. Из второго уравнения системы видим, что $x = 67$. Так как $x = 67$, можно $x$ во первом уравнении заменить на число $67$. Получим уравнение с одной неизвестной, которое мы уже можем решить:
$$
500\cdot 67 + 100y = 22700
$$



%\part{Модули}

%\part{Корни}

\part{Квадратные уравнения}

\chapter{Квадратные уравнения. Начало.}

\section{Приводящие к квадратным уравнениям задачи}

Квадратные уравнения приходят к нам из совершенно разных областей. Их можно встретить во многих задачах геометрии, физики, и теории чисел. 

\subsection{Задача о стороне квадрата}
Простейшая задача, которая приходит к нам из геометрии, формулируется так:

\begin{problem}
Чему равна сторона квадрата с площадью 9 $\text{см}^2$?
\end{problem}
~\\ \textit{Решение:} ~\\ Пусть $x$ --- сторона этого квадрата. Известно, что площадь квадрата --- квадрат его стороны. Используя этот факт, составляем уравнение: $$x^2 = 9.$$ Извлекаем корень из обоих частей, получаем уравнение с модулем: $$|x| = 3.$$ 
Отсюда получаем два решения: $-3$ и $3$. Разумеется, в условиях данной задачи отрицательное решение нам не нужно, получаем ответ: $x=3$ см.

\subsection{Задача о формате бумаги}

Многие из нас пользуются бумагой разных форматов. Пожалуй, самый известный из них --- это формат А4 на рис.\,\ref{fig:A4}. 

\begin{figure}[h]
\label{fig:A4}
\centering
\begin{tikzpicture}
\draw (0,0) rectangle (1.4,2);
\draw (0.7,1) node {A4};
\draw [<->] (1.6,0) -- (1.6,2);
\draw [<->] (0,-0.2) -- (1.4,-0.2);
\draw (0.7,-0.5) node {\tiny $210$ mm};
\draw (1.8,1) node[rotate=90] {\tiny $297$ mm};

\draw (3,0) rectangle (5,2.8);
\draw (4,1.4) node {\huge A3};
\draw [<->] (5.2,0) -- (5.2,2.8);
\draw [<->] (3,-0.2) -- (5,-0.2);
\draw (4,-0.5) node {\tiny $297$ mm};
\draw (5.5,1.4) node[rotate=90] {\tiny $420$ mm};


\end{tikzpicture}
  \caption{Листы бумаги формата А4 и A3}
\end{figure}

\begin{figure}[h]
\label{fig:A4A3}
\centering
\begin{tikzpicture}
\draw (0,0) rectangle (4,2.8);
\draw (1,1.4) node {A4};
\draw (3,1.4) node {A4};
\draw (3.5,0.4) node[color = gray!40] {\Huge A3};
\draw [dotted] (2,0) -- (2,2.8);

\end{tikzpicture}
  \caption{Лист бумаги формата А3 и его разбиение на два листа A4}
\end{figure}

Некоторые также знают, что если склеить по длине два листа А4, мы получим большой лист формата А3 (тоже на рис.\,\ref{fig:A4}). Если же мы решим сложить лист А4 пополам, мы получим лист формата А5. Аналогично для форматов начиная от огромного А0 (его площадь --- 1 квадратный метр) до миниатюрного A8.

Очень примечательная особенность семейства форматов А состоит в том, что \textit{соотношение длины к ширине должно быть одинаковым для любого номера}. 

Хороший вопрос: \textit{а зачем это нужно?} На самом деле, на это есть легкий ответ. Предположим, что вы нарисовали некоторый рисунок на страничке формата А4, a решили перепечатать его в формате А3, чтобы повесить его себе на стену. Так как соотношение длины и ширины страницы не меняется, ваш рисунок можно отобразить <<в масштабе>>, то есть при переходе между форматами ваш рисунок просто увеличится, то есть:

\begin{itemize}
 \item Рисунок не растянется в одном из направлений;
 \item Края рисунка не будут обрезаны.
 \end{itemize} 

Возникает следующая задача:

\begin{problem}
Пусть $l$ -- длина большого листа, $w$ -- его ширина. Каким должно быть соотношение длины к ширине листа $$\frac{l}{w},$$ чтобы сложив его пополам по длине листа, соотношение новой длины $w$ и ширины $\frac{l}{2}$ не изменилось?
\end{problem}
Попробуйте решить эту задачу самостоятельно; если запутаетесь, можно подглянуть в решение.
~\\ \textit{Решение:} ~\\
Соотношение не должно измениться --- значит должно иметь местo равенство:
$$\frac{l}{w} = \frac{w}{l/2}.$$
Это то же самое, что
$$\frac{l}{w} = 2\frac{w}{l}.$$
Нам нужно узнать, чему равно $\frac{l}{w}$. Обозначим эту дробь за неизвестное $x$ и перепишем уравнение: $$x = \frac{2}{x}$$ 
Дальше домножаем обе части на $x$ и получаем уравнение
$$x^2 = 2$$
Отсюда $x =\pm \sqrt{2}$, ввиду положительности ответ $x = \sqrt{2}$. Вспоминаем, что скрывалось за неизвестной $x$ и получаем, что требуемое соотношение длины и ширины листа в таком случае следующее:
$$\frac{l}{w} = \sqrt{2} \approx 1.4$$

Таким образом, если мы хотим, чтобы при складывании некоторого листочка вдвое соотношения сторон не менялись, надо чтобы его длина была в $\sqrt{2}$ раз больше, чем его ширина.

%TODO Задачки с бумагой

\subsection{Бросок тела вверх}

До этого момента нам встречались уравнения, которые требовали только извлечения корня. Ясно, что все подобные задачи решаются достаточно просто. Далее мы рассмотрим задачу, которая приводит к \textit{чуть более} сложному квадратному уравнению.

\begin{problem}

Тело брошено вертикально вверх с поверхности земли с начальной скоростью $30$ м/с. Ускорение свободного падения $g$ считать равным $10$. Через сколько времени тело упадёт на землю?

\end{problem}
~\\ \textit{Решение:} ~\\
Из физики известно, что высота тела, брошенногo с начальной скоростью $v_0$, вычисляется по формуле $$v_0t - \frac{gt^2}{2},$$ где $g$ --- ускорение свободного падения, $t$ --- время, прошедшее с момента броска.

Нас спрашивают: через сколько времени тело упадёт на землю? Значит высота тела должна быть $0$ м. Так и запишем: $$v_0t - \frac{gt^2}{2} = 0.$$ Теперь подставим известные данные --- начальную скорость и ускорение свободного падения: 
$$30t - \frac{10t^2}{2} = 0.$$
Получим уравнение:
$$30t - {5t^2} = 0.$$
С такими уравнениями мы ещё не встречались. Напомним, время $t$ здесь выполняет роль неизвестной величины. 

Здесь не только неизвестное в квадрате ($t^2$), но и неизвестное в первой степени ($t$) участвует. Как же его решать? Подумайте, как можно переписать уравнение так, чтобы его стало легко решить.

На самом деле здесь всё просто: давайте вынесем $t$ за скобки и посмотрим, что получится: $$t(30 - 5t) = 0 .$$ Итак, произведение двух чисел равно нулю тогда, когда хотя бы один из множителей равен нулю. Поэтому получаем два уравнения: $$t = 0, \qquad \qquad 30 - 5t = 0.$$
Когда $t = 0$, по условию задачи тело действительно ещё на поверхности земли, но нам это не особо интересно. С лёгкостью решаем второе уравнение:
$$5t = 30$$
Отсюда $t = 6$ секунд.

\begin{exercise}
\begin{enumerate}
\item Тело брошено вертикально вверх с поверхности Луны с начальной скоростью $20$ м/с. Ускорение свободного падения на Луне $g$ считать равным $1,5$ м/с$^2$. Через сколько времени тело упадёт на поверхость спутника?
\item Тело брошено вертикально вверх с поверхности Юпитера с начальной скоростью $30$ м/с. Ускорение свободного падения $g$ считать равным $24$ м/с$^2$. Через сколько времени тело упадёт на поверхность планеты?
\end{enumerate}
\end{exercise}

Как видим, с помощью такой несложной техники, как вынесение за скобки, мы можем решать задачи такого большого масштаба! К сожалению, все квадратные уравнения так не решить, в этом мы убедимся в следующем пункте.

\section{Выделение полного квадрата}

%TODO

\section{Немного про комплексные числа}

%TODO

\section{Некоторые задачи}

\begin{problem}
Выделить полный квадрат у следующих выражений:
\begin{enumerate}
\item $x^2 + 2x + 10$
\item $x^2 + 2x + \sqrt 2$
\item $4x^2 + 2x + 1$
\item $9x^2 + 6xy + 3y^2$
\end{enumerate}
\end{problem}

\begin{problem}
Решить уравнения, выделив полный квадрат:
\begin{enumerate}
\item $x^2 + 2x + 10 = 0$ (комплексные корни!)
\item $2x^2 + 4x - 5 = 0$
\item $9x^2 - 3x - 6 = 0$
\end{enumerate}
\end{problem}

\begin{problem}
Выяснить, при каких значениях меняющегося параметра $a$ уравнение
$x^2 + 2x + a = 0$ имеет
\begin{itemize}
\item Два вещественных корня;
\item Один вещественный корень;
\item Имеет только комплексные корни.
\end{itemize}
\end{problem}

\begin{problem}
Решить текстовые задачи:
\begin{enumerate}
\item Произведение двух натуральных чисел, одно из которых на $1$ больше другого, равно $210$. Найдите эти числа. 
\item Ширина прямоугольника – $x$, длина на $4$ см больше ширины, а площадь равна $60$. Найдите стороны прямоугольника. 
\end{enumerate}
\end{problem}

\section{Общее решение квадратного уравнения}
\label{sec:quad_gen}

Решаем уравнение $$ax^2+bx+c = 0.$$
Выделим полный квадрат у выражения $ax^2+bx+c$:
\begin{align*}
ax^2+bx+c = 
ax^2 + 2(\sqrt{a}x)\left(\frac{1}{2}b\frac{1}{\sqrt a}\right) + c = 
\left(ax^2 + 2\left(\sqrt{a}x\right)\left(\frac{1}{2}b\frac{1}{\sqrt a}\right) + \frac{b^2}{4a}\right) - \frac{b^2}{4a} + c = \\ = \left(\sqrt{a}x + \frac{b}{2\sqrt{a}}\right)^2 - \frac{b^2}{4a} + c .
\end{align*}

Вернёмся к уравнению $ax^2+bx+c = 0$. Оно эквивалентно следующей цепочке уравнений:

\begin{align*}
\left(\sqrt{a}x + \frac{b}{2\sqrt{a}}\right)^2 - \frac{b^2}{4a} + c &= 0 && \text{выделение полного квадрата} \\ 
\left(\sqrt{a}x + \frac{b}{2\sqrt{a}}\right)^2 &=  \frac{b^2}{4a} - c && 
\text{перенос в правую часть} \\ 
\left(\sqrt{a}x + \frac{b}{2\sqrt{a}}\right) &=  \pm \sqrt{\frac{b^2}{4a} - c} && 
\text{извлечение квадратного корня} \\ 
\sqrt{a}x + \frac{b}{2\sqrt{a}} &=  \pm \sqrt{\frac{b^2 - 4ac}{4a}} && 
\text{преобразование под корнем} \\ 
\sqrt{a}x &=  \pm {\frac{\sqrt{b^2 - 4ac}}{\sqrt{4a}}} - \frac{b}{2\sqrt{a}} && 
\text{перенос и преобразование дроби} \\ 
x &=  \pm {\frac{\sqrt{b^2 - 4ac}}{2a}} - \frac{b}{2a} && 
\text{деление на $\sqrt a$, взятие корня} \\ 
x &=  \frac{-b \pm {\sqrt{b^2 - 4ac}}}{2a} && 
\text{преобразование к классическому виду} 
\end{align*}

Отметим, что чаще всего формулу для корней записывают так, говоря, что $\boxed{D = b^2 - 4ac}$ есть дискриминант квадратного уравнения:
$$\boxed{x = \frac{-b \pm \sqrt D}{2a}}.$$
Разумеется, в преобразованиях предполагалось, что $a \neq 0$. 

\begin{problem}
\textit{Изучение общей формулы. }
\begin{enumerate}
\item При каких $a$, $b$, $c$ уравнение имеет два, один, или не имеет вещественных корней?
\item Придумать уравнение с корнями: $x_1 = 1$ и $x_2 =3$.
\end{enumerate}
\end{problem}


\begin{problem}
С помощью полученной формулы решить уравнения:

\begin{enumerate}
\item $2x^2 + 3x  + 1 = 0$
\item $x^2 - 7x + 10 = 0$
\end{enumerate}

\end{problem}

\section{Предложение о корнях}

\begin{proposition}
\label{prop: roots}
Уравнение $(x-x_1)(x-x_2)=0$ имеет корни $x_1$ и $x_2$.
\end{proposition}

\begin{proof}
Проверяется подстановкой корней $x_1$ и $x_2$ в уравнение $(x-x_1)(x-x_2)=0$.
\end{proof}

Используем это утверждение для решения нескольких простых задач.

\begin{problem}
Какие корни имеет уравнение $$(x-3)(x+7\sqrt 5) = 0 \,?$$
\end{problem}
\begin{problem}
Предложить квадратное уравнение, которое имеет корни $42$ и $\pi$ \footnote{предложить второе такое уравнение}
\end{problem}
\begin{problem}
Применить Предложение \ref{prop: roots} к решению уравнения $$x^2 + 6x + 9 = 0$$
\end{problem}


\section{Теорема Виета}

Рассмотрим \textbf{приведённое}\footnote{$a = 1$ в формуле $ax^2+bx+c = 0$} квадратное уравнение $$x^2 + px + q = 0.$$
Нам уже известнa общая формула для квадратного уравнения: 
$$ x =  \frac{-b \pm {\sqrt{b^2 - 4ac}}}{2a} $$
В данном случае $a = 1$, $b = p$, $c = q$, то есть формула для приведенного уравнения принимает вид:
$$ x = \frac{-p \pm {\sqrt{p^2 - 4q}}}{2} $$
Найдём сумму корней $x_1$ и $x_2$ квадратного уравнения $x^2 + px + q = 0$:
$$ x_1 + x_2 =  \frac{-b + {\sqrt{b^2 - 4c}}}{2} + \frac{-b - {\sqrt{b^2 - 4c}}}{2} = -\frac{b}{2} - \frac{b}{2} = -b. $$
Также найдём произведение корней квадратного уравнения:
\begin{align*}
 x_1 \cdot x_2 =  \frac{-b + {\sqrt{b^2 - 4c}}}{2} \cdot \frac{-b - {\sqrt{b^2 - 4c}}}{2} = \frac{(-b + {\sqrt{b^2 - 4c})(-b - {\sqrt{b^2 - 4c})}}}{4} = \\
= \frac{b^2-(b^2 - 4c)}{4} = \frac{4c}{4} = c
\end{align*}
Полученный результат можно привести как теорему:
\begin{theorem}[Теорема Виета]
Для корней $x_1$ и $x_2$ \textbf{приведённого} квадратного уравнения
$$x^2 + px + q = 0$$ выполняются следующие равенства:
$$ \boxed{x_1 + x_2 = -p \qquad x_1 \cdot x_2 = q}$$
\end{theorem}
Данную теорему можно использовать для решения многих квадратных уравнений. Применим её в решении задач, но сначала приведём пример её использования.

\begin{example}[Применение теоремы Виета]
Решить уравнение $$x^2 + 3x -4 = 0.$$ 
\end{example}
\textit{Решение:} \\ По теореме Виета известно, что сумма корней данного уравнения есть число $-p = - 3$. В то же время, их произведение должно быть равно $q = -4$. 

Таким образом, нам надо подобрать такие два числа, сумма которых равна $-3$, а их произведение равно $-4$. Это несложно: $-4$ и $1$ это такие два числа. По теореме Виета это и будут корни уравнения $x^2 + 3x -4 = 0.$

\begin{problem}
С помощью теоремы Виета решить уравнения:
\begin{enumerate}
\item $x^2 - 7x + 10 = 0$;
\item $x^2 - 4x + 21 = 0$;
\item $x^2 - 13x + 42 = 0$;
\item $x^2 + 2021x + 2022 = 0$;\footnote{попробовать решить с помощью общей формулы ;)}
\item $x^2 - 200x - 404 = 0$.\footnote{аналогично ;)}
\end{enumerate}
\end{problem}


Дальше последует серия задач для тех, кто не верует в силу теоремы Виета:


\begin{problem}
С помощью теоремы Виета решить уравнения:
\begin{enumerate}
\item $2x^2 - 14x + 20 = 0$;
\item $3x^2 - 12x + 63 = 0$;
\item $-x^2 + 13x - 42 = 0$;
\end{enumerate}

\begin{enumerate}
\item Какова особенность этих уравнений по сравнению с предыдущей задачей? \item Удовлетворяют ли уравнения теореме Виета? Почему нет?
\item Можно ли привести их к нужному виду и использовать теорему Виета?
\end{enumerate}

\end{problem}

Выходит, что теорема Виета применима и для неприведённых уравнений! Уточним наши выводы с помощью следующего замечания:

\begin{remark}
\label{com:viete}
Несложно \textbf{проверить самостоятельно},\footnote{проведите рассуждения, аналогичные ситуации с приведённым квадратным уравнением: посчитайте сумму и произведение корней квадратного уравнения $ax^2 +bx+c = 0$ используя общую формулу, полученную на странице \pageref{sec:quad_gen}} что для уравнения в неприведённой форме $$ax^2 +bx+c = 0$$ для корней $x_1$ и $x_2$ верны равенства 
$$ \boxed{x_1 + x_2 = -\frac{b}{a} \qquad x_1 \cdot x_2 = \frac{c}{a}}\,.$$
\end{remark}


\chapter{Частные случаи}

В предыдущих разделах мы рассмотрели различные методы решения квадратных уравнений. Здесь мы подведём некоторый итог, дадим некоторые рекомендации практического характера, и приведём сооответствующие примеры. Это поможет закрепить знания и ускорить процесс решения уравнений.

На частные случаи всегда полезно обратить внимание, так как подобные уравнения часто встречаются и их нужно быстро решать.

\section{Квадратное уравнение <<под ключ>>}
~\\~Квадратное уравнение вида $$\boxed{(x - s)(x - t)= 0}$$ ($s$ и $t$ --- некоторые числа) не надо решать --- оно уже решено! Легко проверить подстановкой, что $s$ и $t$ --- его корни.\footnote{да, это мы повторили предложение \ref{prop: roots}}


\section{Квадратное уравнение без $x$ в первой степени}
~\\~Уравнение вида $$\boxed{ax^2 + c= 0}$$ решается элементарно переносом $c$ в правую часть, а затем взятием корня с обеих сторон равенства.
\\ \textit{Замечание:}
не забывайте, что при взятии корня знаки могут быть разными. Например, у уравнения $$x^2 - 1 = 0$$ два корня: $1$ и $-1$.\footnote{убедитесь в этом}


\section{Квадратное уравнение без свободного члена}
~\\~Уравнение вида $$\boxed{ax^2 + bx = 0}$$ решается вынесением $x$ за скобки. Сразу приведём пример.

\begin{example}
Решить уравнение $\frac{3}{2}x^2 - 6x = 0$.
\end{example}
\textit{Решение:} \\
\begin{align*}
\frac{3}{2}x^2 - 6x &= 0 \\
x\left(\frac{3}{2}x - 6\right) &= 0 \\ 
\boxed{x=0} \text{\quad или \quad} &\frac{3}{2}x - 6 = 0 \\
&\qquad\frac{3}{2}x = 6 \\
&\qquad \boxed{x = 4}.
\end{align*}

\section{Скрытая формула сокращённого умножения}
~\\~Не следует бросаться на амбразуру и \textit{с помощью формулы} решать уравнение вида $$\boxed{a^2x^2 + 2abx + b^2 = 0}\,.$$ Следует заметить, что в левой части скрыта формула сокращённого умножения --- квадрат суммы. Используя этот факт, перепишем уравнение: $$(ax+b)^2 = 0.$$ Ответ следует немедленно: $$x = -\frac{b}{a}.$$
Лучше просто привести пример:
\begin{example}
Решить уравнение $4x^2 + 4x + 1 = 0$
\end{example}
\textit{Решение:} \\
Видим, что в левой части применима формула квадрата суммы. Применяя её, получаем уравнение
$$(2x + 1)^2 = 0.$$ Отсюда легко найти корень: $$x = -\frac{1}{2}.$$ Конечно же, аналогичная ситуация наблюдается и когда в левой части уравнения можно увидеть формулу квадрата разности.

\section{Магия коэффициентов}
Начнём с примера:
\begin{example}
Решить уравнение $6x^2-3x-3 = 0$
\end{example}
\textit{Решение:} \\
Можно заметить, что сумма коэффициентов данного квадратного уравнения ($6$, $-3$, $-3$) равна нулю. 
\begin{quote}
--- Ну и что?
\end{quote}
Дело в том, что если вместо $x$ в уравнение подставить $1$, мы и получим сумму этих трёх коэффициентов в левой части: $$6-3-3 = 0.$$ А ведь это значит, что мы нашли один из корней, ведь при подстановке $x=1$ уравнение превращется в верное тождество. 
\begin{quote}
--- Окей, мы нашли один корень: $x_1 = 1$. И что дальше?
\end{quote}
А дальше мы вспомним, что на стр.~\pageref{com:viete} мы указали на то, что произведение корней квадратного уравнения $x_1\cdot x_2 = \frac{c}{a}$. Зная, что $x_1 = 1$, коэффициент $c = -3$, получим, что $x_2 = -\frac{1}{2}.$

Можно сделать общий вывод:
\begin{quote}
\textbf{Если сумма коэффициентов квадратного уравнения $ax^2+bx+c = 0$ равна нулю, то есть $\boxed{a+b+c = 0}$, то его корни имеют вид $$x_1 = 1, \qquad x_2 = \frac{c}{a}$$}
\end{quote}

Точно таким же образом можно показать, что
\begin{quote}
\textbf{Если для коэффициентов квадратного уравнения $ax^2+bx+c = 0$ выполняется $$\boxed{a+c = b}\,,$$ то его корни имеют вид $$x_1 = -1, \qquad x_2 = -\frac{c}{a}.$$}
\end{quote}

\begin{remark}
Несмотря на то, что такие ситуации с коэффициентами могут показаться выдуманными, уравнения такого вида очень часто встречаются на экзаменах и тестах, так как создателю теста гораздо проще придумать такое уравнение с красивыми корнями. Автор этого текста не стесняясь очень часто сам предлагает к решению такие уравнения. 
\end{remark}

\begin{exercise}
Придумать по два уравнения для каждого из приведённых частных случаев (их здесь $6$) и решить каждое из них.
\end{exercise}

\begin{exercise}
Определить тип частного случая, решить уравнения этим способом:
\begin{enumerate}
\item $3x^2 + x - 4 = 0$
\item $(x-2)(x+6\sqrt{7}) = 0$
\item $36x^2 + 24x + 2 = 0$
\item $122x^2 - 366x = 0$
\item $2x^2+7x+5 = 0$
\end{enumerate}
\end{exercise}

\part{Функции}
\chapter{Мотивировка и определение функции}
\section{Мотивировка 1. Площадь квадрата.}
Представим себе следующую ситуацию: мы хотим узнать площадь квадрата, зная его сторону $x$, где $x$ --- некоторое число ($3$, $\sqrt 7$, $\pi$,\ldots). Ясно, что для получения площади квадрата всё, что нужно сделать --- это перемножить число $x$ на себя, то есть посчитать $x^2$.

В данном примере \textit{функцией} мы называем такую операцию над числом $x$, которая в результате даёт квадрат числа $x$. Функцию (с обозначением $f$) получения площади квадрата можно также выразить на языке символов:
$$
x \xmapsto{f} x^2.
$$
Данная запись читается так:
\begin{quote}
Функция $f$ переводит число $x$ в его квадрат.
\end{quote}
Эта запись позволяет прочувствовать определение функции как некоторого механизма, который принимает некоторое число и возвращает какой-то новый объект.
Несколько примеров работы функции $f$ возведения в квадрат:
\begin{align*}
1 &\xmapsto{f} 1 \\
2 &\xmapsto{f} 4 \\
\sqrt 2 &\xmapsto{f} 2 \\
\pi &\xmapsto{f} \pi^2.
\end{align*}

Классический способ показать, что именно делает функция --- это записать такое выражение: $$f(x) = x^2.$$
$f(x)$ читается как <<$f$ от $x$>> или <<функция $f$ от $x$>>.

\begin{remark}
Может возникнуть логичный вопрос: <<Зачем вводить новое и никому не нужное обозначение $f(x)$?>> Дело в том, что функций в одной задаче может быть много.
Как тогда различить $$f(x) = x^2 \qquad \text{от} \qquad f(x) = x + 3 \text{ ?}$$
Если же задать $g(x) = x+3$, тогда всем будет ясно, что $f(3) = 9$, a $g(3) = 6$.
\end{remark}

\begin{example}
Чему равно $f(g(3))$?

Следует запомнить, что функции <<распутываются>> изнутри:вполне разумно посчитать то, что внутри: $g(3)$, а затем уже полученное значение подставить в функцию $f$.

\end{example}

\begin{exercise}
Пусть заданы следующие функции:
\begin{itemize}
\item $f(x) = x^2 + 3x$
\item $g(x) = \frac{2}{x}$
\item $h(x) = 2^x$
\end{itemize}

Вычислите:
\begin{enumerate}
\item $f(f(1))$
\item $g(g(2))$
\item $g(f(b))$
\item $h(g(f(a)))$
\end{enumerate}
\end{exercise}

\begin{exercise}
Ниже записано, какие числа получаются при использовании функции $h$ для преобразования некоторых чисел.
\begin{align*}
1 &\xmapsto{h} 1 \\
2 &\xmapsto{h} 8 \\
3 &\xmapsto{h} 27 \\
10 &\xmapsto{h} 1000.
\end{align*}
Что это может быть за функция? Запишите ответ с помощью записи со стрелками и в виде $h(x) = \ldots$
\end{exercise}

\begin{exercise}
Ниже записано, какие числа получаются при использовании функции $q$ для преобразования некоторых чисел.
\begin{align*}
2 &\xmapsto{q} 5 \\
3 &\xmapsto{q} 10 \\
4 &\xmapsto{q} 17 \\
5 &\xmapsto{q} 26.
\end{align*}
Что это может быть за функция? Запишите ответ с помощью записи со стрелками и в виде $q(x) = \ldots$
\end{exercise}

\section{Мотивировка 2. Пройденное расстояние.}

Большое число интересных примеров функций приходит к нам из области физики. Здесь расскажем о простейшем примере: расстояние, пройденное за некоторое время.

Предположим, что машина в момент времени $0$ часов движется с постоянной скоростью $80$ км/ч. Тогда пройденное расстояние $S$ за $t$ часов будет равно $80\cdot t$. 

Mожно рассмотреть расстояние $S$ как функцию от \textit{переменной}\footnote{в примере с площадью квадрата в качестве переменной мы рассматривали $x$} $t$:
$$ S(t) = 80t. $$
Получается, что зная среднюю скорость движения по автодороге и имея при себе наручные часы, можно примерно знать пройденное расстояние в любой момент времени.

\begin{exercise}
\begin{enumerate}
\item Предложите функцию $P(t)$ от времени, которая поможет рассчитать оставшееся расстояние до города, который в начале пути автомобиля находится на расстоянии $880$ километров. Скорость машины считать постоянной --- $80$ км/ч.
\item С помощью функции $P(t)$ найти $P(2)$, $P(5)$; объяснить смысл данных значений.
\end{enumerate}


Оставшееся расстояние есть длина всего пути без пройденного расстояния.

\end{exercise}

\begin{remark}
На самом деле, функции бывают очень разные. Например, можно рассмотреть такую функцию:\footnote{функцию Дирихле}
$$
f(x) = \begin{cases}
1, \text{ если х --- рациональное число} \\
0, \text{ иначе} \\
\end{cases}
$$
Данная запись означает, что если $x$ --- рациональное число,\footnote{представимое в виде дроби, где числитель и знаменатель --- целые числа} то $f(x)$ будет равно $1$, а если нет --- то $f(x) = 0$. Поэтому $$f(1) = 1, \qquad f(0) = 1, \qquad f\left(\frac{35}{46}\right) = 1,$$
но при этом
$$f(\sqrt{2}) = 0, \qquad f(\pi) = 0, \qquad f\left(\frac{\sqrt{2}}{2}\right) = 0.$$

Также часто можно встретить константные функции, такие как $$f(x) = 0, \qquad g(x)=3.14,\qquad h(x) = 425.$$ Заметим, что здесь нет $x$ в правой части выражения. Это означает, что функция не меняется при изменении $x$. На нескольких примерах это означает следующее: $$f(1) = 0,\, f(2) = 0,\, f(100) = 0,\, g(0) = 3.14,\,g(4) = 3.14,\, h(0) = 425,\, h(67) = 425, \ldots$$

Точно так же никого не должны смущать и такие функции: $$f(x) = 4ax,\qquad g(x) = 4x^w.$$ Когда увидите такую запись, следует уточнить, что такое $a$ и $w$ в данных функциях --- какие-то фиксированные числа или произвольные. И уже после уточнения решать ту или иную задачу.

\end{remark}

\section{Область определения и область значений функции}

Мы до сих пор не задались вопросом о том, какие ограничения могут существовать на $x$ --- значения переменной (\textit{аргумента}) и $f(x)$ --- значения функции.

\paragraph{Ограничения на $x$: область определения функции}
Сначала разберёмся с тем, какие значения может принимать аргумент $x$. Откуда вообще могут взяться какие-то ограничения на $x$? Например, посмотрим на такую функцию $$f(x) = 3x^2 + 2x - 1.$$ Нетрудно понять, что какое бы число $x$ мы не подставили, $f(x)$ можно вычислить. В таком случае мы говорим, что функция $f(x)$ имеет область определения $(-\infty, +\infty)$, то есть \textit{определена на всей числовой оси}. 

Теперь рассмотрим такую функцию: $$g(x) = \frac{1}{x}.$$ Вопрос остаётся: есть ли какие-то ограничения на $x$? Ответ должен быть понятен всем: $x$ не может быть нулём, так как на ноль делить нельзя. Любые другие числа, разумеется, подходят, то есть область определения функции $g(x)$ есть $(-\infty, 0)\cup(0, +\infty)$.

В качестве примеров дадим несколько задач:

\begin{exercise}
Укажите области определения следующих функций 
\begin{enumerate}
\item $f(x) = \frac{x}{1+x}$
\item $g(x) = \sqrt x + \frac{1}{x}$
\item $h(x) = \frac{1}{\sqrt{x}}$
\item $\alpha(x) = \frac{1}{x\sqrt{x}}$ 
\item $\beta(x) = \frac{x}{x^2 + x}$ 
\item $\gamma(x) = \frac{x-2}{x^2 - 8x + 6}$
\end{enumerate}
\end{exercise}

\subsection{Ограничения на $f(x)$: область значений функции}
А теперь зададимся таким вопросом: какие значения может принимать сама функция $f(x)$? Например, возьмём ужe упомянутый пример $$f(x) = \frac{1}{x}.$$ Попробуйте придумать такое число $a$, чтобы $f(a) = 0$. 

Если немного поразмышлять, станет понятно, что такого числа найти не получится! Ясно, что для, например, $f(a) = 2$ или $f(a) = \frac{1}{2}$, такое число $a$ найти можно.\footnote{точно?} Да и вообще, должно быть понятно, что эта функция принимает все возможные значения, кроме нуля.\footnote{действительно понятно?} Следовательно, область значений функции $f(x)$ будет состоять из всех чисел, кроме нуля. 

Вообще, вычисление области значений функции иногда очевидно, а иногда заведомо сложнее области определения функции из-за того, что для её нахождения надо перебрать все возможные $x$ из области определения. Поэтому мы коснёмся только основных деталей.

Приведём несколько примеров вычисления области значений функции:

\begin{example}
Укажите области значений следующих функций 
\begin{enumerate}
\item $f(x) = \sqrt{x}$ \\
Нам известно, что \textit{область определения} функции $f(x)$ есть $[0, +\infty)$. Для того, чтобы узнать область значений функции, \emph{надо перебрать все возможные $x$ из области определения}. Начнём с $x = 0$ и закончим большими значениями $x$, которые стремятся к бесконечности. Итак, прикинем несколько значений $f(0) = 0,\, f(1) = 1,\, f(3) = \sqrt{3},\, f(10000) = 100$ и так далее. Видим, что функция корня возрастает и никаких ограничений сверху мы не найдём, поэтому функция принимает значения $$[0, +\infty).$$
\item $g(x) = x$ \\
Здесь область определения и область значений совпадают и составляют всю числовую ось от $-\infty$ до $+\infty$.
\item $g(x) = \frac{1}{x} + 1$ \\
Здесь мы используем то, что уже посчитали в самом начале для функции $f(x) = \frac{1}{x}$. Там область значений состояла из всех чисел, кроме $0$. Функция $g(x)$ отличается от $f(x)$ добавлением единицы, поэтому в данном случае область значений смещается на единицу, то есть имеем $$(-\infty, 1)\cup(1, +\infty).$$
\end{enumerate}
\end{example}

\begin{exercise}
Укажите области значений следующих функций 
\begin{enumerate}
\item $f(x) = \frac{1}{x} - 2$
\item $g(x) = \sqrt{x} - \frac{1}{2}$
\item $h(x) = \sqrt{x+1}$
\end{enumerate}
\end{exercise}


\section{Определение функции (*)}

Предполагается, что на данный момент должно было сложиться представление о таком математическом объекте, как функция. Дадим общее определение функции, которое используется в математике и уточним его смысл.

\begin{definition}
Функцией $f(x)$ называется такое отображение из области определения в область значений, у которого каждому из чисел из области определения соответствует единственное число в области значений.
\end{definition}

Данное определение, безусловно, требует объяснений. Тот факт, что функция есть отображение из области определения в область значений должен быть понятен всем, кто прочитал предыдущие пункты. Теперь, обсудим часть с <<единственностью>>. Определение говорит, что какое бы мы не взяли число $x$, значение $f(x)$ будет единственным. Другими словами, \textit{такая} схема соответствует определению функции:
\begin{align*}
1 &\xmapsto{f} 1 \\
2 &\xmapsto{f} 8 \\
3 &\xmapsto{f} 27 \\
10 &\xmapsto{f} 1000,
\end{align*}
а \textit{такая} ---
\begin{align*}
\boxed{1} &\xmapsto{f} 1 \\
\boxed{1} &\xmapsto{f} 8 \\
3 &\xmapsto{f} 27 \\
10 &\xmapsto{f} 1000 
\end{align*}
--- уже нет, так как в данной схеме мы не можем сказать однозначно, чему равно значение $f(1)$: можно лишь сказать, что оно либо равно $1$, либо $8$. 

Можно спросить --- а зачем такие ограничения? На самом деле, это вполне логичное и практически обоснованное требование.  Приведём пару примеров:

\begin{example}
Например, если $S(t)$ --- некоторая функция пройденного расстояния в момент времени $t$, мы считаем, что в  фиксированный момент времени $t=1$ тело находится на определённом расстоянии от начала пути. Если же $S(1)=100$ и $S(1)=200$ одновременно, то это не совсем хорошо.\footnote{здесь должна быть несмешная шутка про квантовую физику} 
\end{example}

\begin{example}
$f(x) = x^2$ --- действительно функция, так как квадрат любого числа $x$ --- всегда некоторое единственное число. То же самое можно сказать и про все функции, которые мы рассматривали ранее. 
\end{example}

\begin{remark}
Заметим, что \emph{такая} схема не запрещена:
\begin{align*}
-2 &\xmapsto{f} \boxed{4} \\
-1 &\xmapsto{f} \boxed{1} \\
1 &\xmapsto{f} \boxed{1} \\
2 &\xmapsto{f} \boxed{4} 
\end{align*}
Как можно видеть, здесь повторяются значения функции для разных значений $x$. Если посмотреть на определение функции, никаких противоречий мы не обнаружим. Кстати, некоторые наверняка увидели, что по этой схеме может работать функция $f(x) = x^2$.
\end{remark}

\part{Последовательности}

\chapter{Краткое введение в последовательности}

В данной главе мы начнём наше знакомство с одним из самых используемых на практике понятий математики --- числовыми последовательностями, которые, как оказывается, являются частными случаями функций. 

\section{Последовательности и их примеры}
По названию ясно, что последовательности --- это некоторые объекты, которые могут выглядеть вот так:
\begin{align*}
1,\,2,\,3,\,4,\,5; \\
1,\,3,\,5,\,7,\,9; \\
2,\,4,\,6,\,8,\,10; \\
1,\,2,\,3,\,5,\,7,\,11,\,13,\,17.
\end{align*}
Несложно увидеть, что приведённые сверху последовательности подчиняются некоторым законам/правилам (последняя, например, задаёт последовательность, состоящую из первых восьми простых чисел). Эти правила могут быть какими угодно, поэтому такие последовательности тоже встречаются:
\begin{align*}
1,\,0,\,1,\,0,\,1 & \text{\qquad что-то чередующееся;} \\
5,\,5,\,5,\,5,\,5 & \text{\qquad что-то одинаковое;} \\
3,\,1,\,4,\,1,\,5 & \text{\qquad разложение числа пи;} \\
1,\,2,\,3,\,3,\,2,\,1,\,1,\,2,\,3 & \text{\qquad что-то чередующееся.}
\end{align*}

Также последовательности могут быть бесконечными:
\begin{align*}
5,\,5,\,5,\,5,\,5,\ldots & \text{\qquad бесконечный стрик из пятёрок;} \\
3,\,1,\,4,\,1,\,5\ldots & \text{\qquad разложение числа пи;} \\
2^0,\,2^1,\,2^2,\,2^3,\,2^4,\ldots & \text{\qquad степени двойки;} \\
1,\,2,\,3,\,3,\,2,\,1,\,1,\,2,\,3,\ldots & \text{\qquad что-то бесконечно чередующееся.}
\end{align*}

В нашем определении последовательности порядок важен, поэтому последовательности $1,\,2,\,3,\,4$ и $2,\,1,\,4,\,3$ --- разные.

Отдельные числа в последовательностях называются их \textbf{членами}. Как и у функций ($f$, $g$, $h$), у каждой последовательности должно быть своё имя; принято выбирать буквы $a$, $b$, $c$, \ldots. Kаждый её член обозначется именем последовательности с соответствующим индексом внизу справа. Например, для последовательности с именем $a$ $$1,\,2,\,3,\,5,\,7,\,11,\,13,\,17,\,$$ можем указать первые пять членов последовательности 
\begin{align*}
a_1 &= 1 \\
a_2 &= 2 \\
a_3 &= 3 \\
a_4 &= 5 \\
a_5 &= 7 
\end{align*}

\begin{exercise}
Для последовательности $a$ выше, назовите $a_7$ и $a_8$
\end{exercise}

\begin{exercise}
Пусть $a$ -- последовательность всех чётных чисел больше нуля, $b$ -- последовательность всех нечётных чисел больше нуля. Назовите:
\begin{itemize}
\item $a_1,\,a_3,\,a_5$
\item $b_3,\, b_5,\, b_6$
\item $a_3 + b_3$
\item $a_{10}$, $a_{14}$, $a_{20}$, $a_{100}$ (разумеется, не перебирая все элементы)
\item $b_{10}$, $b_{20}$, $b_{100}$
\end{itemize} 
\end{exercise}

\chapter{Формулы последовательностей}

Так же, как и функции, некоторые последовательности могут быть определены (заданы) с помощью формул. 
Вот, например, формула для последовательности положительных чётных чисел: $$  a_i = 2\cdot i$$ Вместо $i$ мы подставляем нужный нам номер по порядку. Отсюда получаем, что, например, $a_{12} = 24$.

\begin{exercise}
Какая формула будет у последовательности положительных нечётных чисел?
\end{exercise}




\part{Комбинаторика}



\begin{problem}
У англичан принято давать ребёнку несколько имён. Сколькими способами в Англии можно назвать ребёнка, если общее число имён равно 300, а ему дают не более трёх имён. Хватит ли этих имён на всех англичан (66 млн.~чел) или же непременно найдутся англичане с одинаковыми именами? (\textit{Виленкин, <<Комбинаторика>>})
\end{problem}

Решение:

\begin{problem}
Можно ли выписать девять чисел $1,2,\ldots,9$ по кругу так, чтобы сумма никаких двух соседних чисел не делилась ни на $3$, ни на $5$, ни на $7$? (\textit{Виленкин, <<Комбинаторика>>})
\end{problem}








%\section{Упражнения}
%Упражнения даны в порядке увеличения сложности, рекомендуется делать их по порядку.
%
%TODO::



%%----------------------------------------------------------------------------------------
%%	PART
%%----------------------------------------------------------------------------------------
%
%\part{Part One}
%
%%----------------------------------------------------------------------------------------
%%	CHAPTER 1
%%----------------------------------------------------------------------------------------
%
%\chapterimage{chapter_head_2.pdf} % Chapter heading image
%
%\chapter{Text Chapter}
%
%\section{Paragraphs of Text}\index{Paragraphs of Text}
%
%\lipsum[1-7] % Dummy text
%
%%------------------------------------------------
%
%\section{Citation}\index{Citation}
%
%This statement requires citation \cite{article_key}; this one is more specific \cite[162]{book_key}.
%
%%------------------------------------------------
%
%\section{Lists}\index{Lists}
%
%Lists are useful to present information in a concise and/or ordered way\footnote{Footnote example...}.
%
%\subsection{Numbered List}\index{Lists!Numbered List}
%
%\begin{enumerate}
%\item The first item
%\item The second item
%\item The third item
%\end{enumerate}
%
%\subsection{Bullet Points}\index{Lists!Bullet Points}
%
%\begin{itemize}
%\item The first item
%\item The second item
%\item The third item
%\end{itemize}
%
%\subsection{Descriptions and Definitions}\index{Lists!Descriptions and Definitions}
%
%\begin{description}
%\item[Name] Description
%\item[Word] Definition
%\item[Comment] Elaboration
%\end{description}
%
%%----------------------------------------------------------------------------------------
%%	CHAPTER 2
%%----------------------------------------------------------------------------------------
%
%\chapter{In-text Elements}
%
%\section{Theorems}\index{Theorems}
%
%This is an example of theorems.
%
%\subsection{Several equations}\index{Theorems!Several Equations}
%This is a theorem consisting of several equations.
%
%\begin{theorem}[Name of the theorem]
%In $E=\mathbb{R}^n$ all norms are equivalent. It has the properties:
%\begin{align}
%& \big| ||\mathbf{x}|| - ||\mathbf{y}|| \big|\leq || \mathbf{x}- \mathbf{y}||\\
%&  ||\sum_{i=1}^n\mathbf{x}_i||\leq \sum_{i=1}^n||\mathbf{x}_i||\quad\text{where $n$ is a finite integer}
%\end{align}
%\end{theorem}
%
%\subsection{Single Line}\index{Theorems!Single Line}
%This is a theorem consisting of just one line.
%
%\begin{theorem}
%A set $\mathcal{D}(G)$ in dense in $L^2(G)$, $|\cdot|_0$. 
%\end{theorem}
%
%%------------------------------------------------
%
%\section{Definitions}\index{Definitions}
%
%This is an example of a definition. A definition could be mathematical or it could define a concept.
%
%\begin{definition}[Definition name]
%Given a vector space $E$, a norm on $E$ is an application, denoted $||\cdot||$, $E$ in $\mathbb{R}^+=[0,+\infty[$ such that:
%\begin{align}
%& ||\mathbf{x}||=0\ \Rightarrow\ \mathbf{x}=\mathbf{0}\\
%& ||\lambda \mathbf{x}||=|\lambda|\cdot ||\mathbf{x}||\\
%& ||\mathbf{x}+\mathbf{y}||\leq ||\mathbf{x}||+||\mathbf{y}||
%\end{align}
%\end{definition}
%
%%------------------------------------------------
%
%\section{Notations}\index{Notations}
%
%\begin{notation}
%Given an open subset $G$ of $\mathbb{R}^n$, the set of functions $\varphi$ are:
%\begin{enumerate}
%\item Bounded support $G$;
%\item Infinitely differentiable;
%\end{enumerate}
%a vector space is denoted by $\mathcal{D}(G)$. 
%\end{notation}
%
%%------------------------------------------------
%
%\section{Remarks}\index{Remarks}
%
%This is an example of a remark.
%
%\begin{remark}
%The concepts presented here are now in conventional employment in mathematics. Vector spaces are taken over the field $\mathbb{K}=\mathbb{R}$, however, established properties are easily extended to $\mathbb{K}=\mathbb{C}$.
%\end{remark}
%
%%------------------------------------------------
%
%\section{Corollaries}\index{Corollaries}
%
%This is an example of a corollary.
%
%\begin{corollary}[Corollary name]
%The concepts presented here are now in conventional employment in mathematics. Vector spaces are taken over the field $\mathbb{K}=\mathbb{R}$, however, established properties are easily extended to $\mathbb{K}=\mathbb{C}$.
%\end{corollary}
%
%%------------------------------------------------
%
%\section{Propositions}\index{Propositions}
%
%This is an example of propositions.
%
%\subsection{Several equations}\index{Propositions!Several Equations}
%
%\begin{proposition}[Proposition name]
%It has the properties:
%\begin{align}
%& \big| ||\mathbf{x}|| - ||\mathbf{y}|| \big|\leq || \mathbf{x}- \mathbf{y}||\\
%&  ||\sum_{i=1}^n\mathbf{x}_i||\leq \sum_{i=1}^n||\mathbf{x}_i||\quad\text{where $n$ is a finite integer}
%\end{align}
%\end{proposition}
%
%\subsection{Single Line}\index{Propositions!Single Line}
%
%\begin{proposition} 
%Let $f,g\in L^2(G)$; if $\forall \varphi\in\mathcal{D}(G)$, $(f,\varphi)_0=(g,\varphi)_0$ then $f = g$. 
%\end{proposition}
%
%%------------------------------------------------
%
%\section{Examples}\index{Examples}
%
%This is an example of examples.
%
%\subsection{Equation and Text}\index{Examples!Equation and Text}
%
%\begin{example}
%Let $G=\{x\in\mathbb{R}^2:|x|<3\}$ and denoted by: $x^0=(1,1)$; consider the function:
%\begin{equation}
%f(x)=\left\{\begin{aligned} & \mathrm{e}^{|x|} & & \text{si $|x-x^0|\leq 1/2$}\\
%& 0 & & \text{si $|x-x^0|> 1/2$}\end{aligned}\right.
%\end{equation}
%The function $f$ has bounded support, we can take $A=\{x\in\mathbb{R}^2:|x-x^0|\leq 1/2+\epsilon\}$ for all $\epsilon\in\intoo{0}{5/2-\sqrt{2}}$.
%\end{example}
%
%\subsection{Paragraph of Text}\index{Examples!Paragraph of Text}
%
%\begin{example}[Example name]
%\lipsum[2]
%\end{example}
%
%%------------------------------------------------
%
%\section{Exercises}\index{Exercises}
%
%This is an example of an exercise.
%
%\begin{exercise}
%This is a good place to ask a question to test learning progress or further cement ideas into students' minds.
%\end{exercise}
%
%%------------------------------------------------
%
%\section{Problems}\index{Problems}
%
%\begin{problem}
%What is the average airspeed velocity of an unladen swallow?
%\end{problem}
%
%%------------------------------------------------
%
%\section{Vocabulary}\index{Vocabulary}
%
%Define a word to improve a students' vocabulary.
%
%\begin{vocabulary}[Word]
%Definition of word.
%\end{vocabulary}
%
%%----------------------------------------------------------------------------------------
%%	PART
%%----------------------------------------------------------------------------------------
%
%\part{Part Two}
%
%%----------------------------------------------------------------------------------------
%%	CHAPTER 3
%%----------------------------------------------------------------------------------------
%
%\chapterimage{chapter_head_1.pdf} % Chapter heading image
%
%\chapter{Presenting Information}
%
%\section{Table}\index{Table}
%
%\begin{table}[h]
%\centering
%\begin{tabular}{l l l}
%\toprule
%\textbf{Treatments} & \textbf{Response 1} & \textbf{Response 2}\\
%\midrule
%Treatment 1 & 0.0003262 & 0.562 \\
%Treatment 2 & 0.0015681 & 0.910 \\
%Treatment 3 & 0.0009271 & 0.296 \\
%\bottomrule
%\end{tabular}
%\caption{Table caption}
%\label{tab:example} % Unique label used for referencing the table in-text
%%\addcontentsline{toc}{table}{Table \ref{tab:example}} % Uncomment to add the table to the table of contents
%\end{table}
%
%Referencing Table \ref{tab:example} in-text automatically.
%
%%------------------------------------------------
%
%\section{Figure}\index{Figure}
%
%\begin{figure}[h]
%\centering\includegraphics[scale=0.5]{placeholder.jpg}
%\caption{Figure caption}
%\label{fig:placeholder} % Unique label used for referencing the figure in-text
%%\addcontentsline{toc}{figure}{Figure \ref{fig:placeholder}} % Uncomment to add the figure to the table of contents
%\end{figure}
%
%Referencing Figure \ref{fig:placeholder} in-text automatically.

%----------------------------------------------------------------------------------------
%	BIBLIOGRAPHY
%----------------------------------------------------------------------------------------

%\chapter*{Bibliography}
%\addcontentsline{toc}{chapter}{\textcolor{ocre}{Bibliography}} % Add a Bibliography heading to the table of contents
%
%%------------------------------------------------
%
%\section*{Articles}
%\addcontentsline{toc}{section}{Articles}
%\printbibliography[heading=bibempty,type=article]
%
%%------------------------------------------------
%
%\section*{Books}
%\addcontentsline{toc}{section}{Books}
%\printbibliography[heading=bibempty,type=book]

%----------------------------------------------------------------------------------------
%	INDEX
%----------------------------------------------------------------------------------------

\cleardoublepage % Make sure the index starts on an odd (right side) page
\phantomsection
\setlength{\columnsep}{0.75cm} % Space between the 2 columns of the index
\addcontentsline{toc}{chapter}{\textcolor{ocre}{Index}} % Add an Index heading to the table of contents
\printindex % Output the index

%----------------------------------------------------------------------------------------

\end{document}
